\mediumtitle{Technology Transfer and Public Engagement}

% \mediumtitle{Outreach activity}

\smallskip

Alessia Visconti's activities include, in addition to the development and release of open-source software (see the \textsc{Development of research software} section), the initiatives listed below.

\smallskip

\begin{doubletablelist}

	\dtlist{Jun 2026 - Present}{\textbf{Selected for Funds TOgether}, the crowdfunding programme of the University of Turin dedicated to high-impact and innovative projects with social and environmental impact. Alessia Visconti (PI of the project) and her team (Paola Berchialla, Barbara Sini, Matteo Bracco) passed a selection process assessing both the quality of the project proposal, including its scientific value and innovativeness, and the suitability of the project objectives for crowdfunding. %The latter included the coherence between the campaign fundraising target and the team potential in terms of digital communication, target community, and fundraising experience, as well as the quality of the content and communication strategy. 
The proposed project aims at understanding how climate change affects psychological wellbeing
%. Indeed, an increasing number of people, particularly young people, report anxiety, stress, insomnia, and fear about the future related to the climate crisis. To address this, the project will 
by developing a citizen-science app to collect longitudinal data on psychological wellbeing and integrate them with climate and environmental data. The aim is to create the first participatory observatory on the relationship between climate and mental health, raising public awareness and supporting concrete interventions at the local level.}

	\dtlist{Sep 2025 - Present}{\textbf{Participation in the \emph{Penne Amiche della Scienza} project}. The project follows the format of \emph{``Letters to a Pre-Scientist"} and aims at introducing primary and lower-secondary school students to the world of science through an exchange of letters with researchers. These letters provide an opportunity to share researchers' educational and research journeys, answer students' questions, and stimulate their interest in scientific disciplines. Alessia Visconti participated in the 2025-26 edition of the programme, which involved sending four letters (one introductory letter and three responses) and visiting the participating school (\emph{Scuola Secondaria I grado Antonio Gramsci}, Class 2A). She has also enrolled in the 2026-27 edition and will be paired with a school in September 2026.}

	\dtlist{Dec 2025 - Mar 2026}{\textbf{Recording of a video lecture} (in Italian) entitled \emph{``Introduction to Large Language Models''} (Module 3) within the transdisciplinary Massive Open Online Course (MOOC): \emph{``Artificial Intelligence in Veterinary Medicine: Perspectives and Applications''}, which was funded by the National Recovery and Resilience Plan (PNRR), Mission 4, Education and Research.}

	\dtlist{Jul 2025 - Present}{\textbf{Establishment of the start-up ``Repertor.IO''}, approved by the Department Council on June 11th, 2024 and established on July 14th, 2025. Repertor.IO was recognised as an academic spin-off by the Board of Directors of the University of Turin on July 18th 2025. Alessia Visconti's role is \emph{Chief Research Communication Manager}.}

	\dtlist{Jul 2017 -- Feb 2023}{\textbf{Member and mentor} of the \emph{Artificial Intelligence Club for Gender Minorities}, which aimed at promoting gender diversity in the artificial intelligence and scientific community via meetups and mentorship. Alessia Visconti organised workshops on collaborative data science using Git and GitHub. From May 2018, she also co-organised the group monthly journal club.}

	\dtlist{Mar 2016 -- Sep 2017}{\textbf{Member and mentor} of the \emph{RLadies London} community and of the \emph{Researc[her] Research} community. These groups aim at promoting gender diversity in the R and STEM community via meetups, mentorship and global collaboration. With \emph{Researc[her]}, Alessia Visconti gave speeches on reproducibility and workflow development.}

\end{doubletablelist}



\mediumtitle{Development of research software}

Alessia Visconti is the sole developer unless otherwise specified. The software is distributed under an open-source licence and is available at \url{https://github.com/alesssia} or upon request.

\smallskip

\begin{singletablelist}


	\stlist{AID-ISA}{biclustering algorithm leveraging \emph{a priori} knowledge~\cite{Vis13a}}

	\stlist{CDoT}{exact algorithm for answering \emph{Maximum a Posteriori} queries on tree-structured probabilistic graphical models~\cite{Esp13} (developed in collaboration with Roberto Esposito)}

	\stlist{famCNV (v2.0)}{algorithm for testing the association between copy-number variation (CNV) and quantitative phenotypes in families (developed in collaboration with Mario Falchi)}

	\stlist{GOClust}{co-clustering algorithm that leverages biological information extracted from the Gene  Ontology to improve the quality of the results~\cite{Vis11c,Cor09b}}

	\stlist{MotifsLinker}{algorithm for extracting protein motifs and grouping them into families using a constrained co-clustering approach~\cite{Cor09a}}

	\stlist{PicNic}{tool for searching topologies within chemical molecules, which has been used to support teaching at the Department of Molecular Biotechnology and Health Sciences, University of Turin}

	\stlist{PopPAnTe}{algorithm for pairwise association testing of quantitative data in related samples~\cite{Vis16}}

	\stlist{RGO}{algorithm for grouping Gene Ontology terms according to their function to obtain a more compact and informative ontology~\cite{Vis11a,Vis10a}}

	\stlist{SPOT}{algorithm for the exhaustive search of frequent motifs in biological sequences~\cite{Vis08}}

	\stlist{YAMP}{reproducible pipeline for the analysis of metagenomic data~\cite{Vis18b}}

\end{singletablelist}